\documentclass[a4paper,11pt]{article}

\usepackage[T1]{fontenc}
\usepackage[utf8]{inputenc}
\usepackage[UKenglish]{babel}
\usepackage{amsmath,amssymb,amsthm}
\usepackage[ruled,vlined,linesnumbered]{algorithm2e}
\usepackage{tikz}
\usetikzlibrary{positioning,shadows}
\usepackage{hyperref}
\usepackage{xspace}
\usepackage{xcolor}
\usepackage[margin=1in]{geometry}
\usepackage{microtype}

% theorem environments
\newtheorem{theorem}{Theorem}[section]
\newtheorem{lemma}[theorem]{Lemma}
\newtheorem{corollary}[theorem]{Corollary}
\newtheorem{proposition}[theorem]{Proposition}
\newtheorem{definition}[theorem]{Definition}
\newtheorem{remark}[theorem]{Remark}

% problem definition style
\tikzstyle{mybox} = [draw=black, fill=white, drop shadow, very thick, rectangle, rounded corners, inner sep=10pt, inner ysep=12pt]
\tikzstyle{fancytitle} =[fill=white, text=black]
\newcommand{\pbdef}[3]{
\noindent
\begin{tikzpicture}
\node [mybox] (box){%
    \begin{minipage}{0.94\textwidth}
        \textit{Input: }{#2}
        \\
        \textit{Task: }{#3}
    \end{minipage}
};
\node[fancytitle, right=10pt] at (box.north west) {#1};
\end{tikzpicture}
}

\newcommand{\br}[1]{\mathopen{}\left( #1 \right)}
\newcommand{\brc}[1]{\mathopen{}\left\{ #1 \right\}}
\newcommand{\spr}[1]{\mathopen{}\left| #1 \right|}
\newcommand{\fl}[1]{\mathopen{}\left\lfloor #1 \right\rfloor}
\newcommand{\cl}[1]{\mathopen{}\left\lceil #1 \right\rceil}
\newcommand{\angl}[1]{\mathopen{}\langle #1 \rangle}
\newcommand{\e}[1]{\exp\left\{ #1 \right\}}
\newcommand{\diam}{\operatorname{diam}}
\newcommand{\NPcomplete}{\textnormal{NP-Complete}\xspace}
\newcommand{\NPhard}{\textnormal{NP-Hard}\xspace}
\newcommand{\Cent}{\texttt{Cent}}
\newcommand{\APP}{\texttt{APP}}
\newcommand{\OPT}{\texttt{OPT}}
\newcommand{\COST}{\texttt{COST}}
\newcommand{\argmin}{\operatorname*{argmin}}
\newcommand{\argmax}{\operatorname*{argmax}}
\newcommand{\cov}{\operatorname{cov}}
\newcommand{\E}[1]{\mathbb{E}\!\left[#1\right]}


\bibliographystyle{plainurl}

\title{Polylogarithmic Approximation for Covering and Connecting Multi-Interface Networks}

\author{Micha\l{} Szyfelbein\thanks{Gda\'nsk University of Technology, Poland. \texttt{michal.szyfelbein@pg.edu.pl}. \href{https://orcid.org/0009-0009-9894-9671}{ORCID: 0009-0009-9894-9671}}
\and
Camille Richer\thanks{Universit\'e Paris-Dauphine, PSL Research University, CNRS, UMR 7243, LAMSADE. \texttt{camille.richer@dauphine.eu}. \href{https://orcid.org/0009-0000-3636-6571}{ORCID: 0009-0000-3636-6571}}}

\date{}

\begin{document}

\maketitle

\begin{abstract}
We study problems related to connecting multi-interface networks of wireless devices. These problems can be modeled using graphs, where vertices represent the devices and edges represent potential communication  links. Each vertex can activate multiple interfaces, and a connection between two vertices is established if they share at least one common active interface.
However activating an interface induces a cost that depends on the type of the interface and on the vertex that activates it. We consider two problems arising in multi-interface networks: \coverage and \connectivity. In the \coverage problem, every connection defined in the network must be established, while in the \connectivity problem, it is only required that established connections form a connected spanning subgraph of the network. The solution should also minimize either the maximum cost incurred by a vertex or the total cost incurred by all vertices. In this work we are interested in approximating the former of the two cost criterions.

We model both problems using Integer Linear Programming (ILP) and we design approximation algorithms based on a randomized rounding of the solution of the linear programming relaxation. For the \coverage problem, this yields an $O(\log n)$-approximation algorithm, where $n$ is the number of vertices. This result is tight, since the problem generalizes \setcover. This improves upon the $O(b\cdot\log n)$-approximation algorithm for the min-sum version of the problem, where $b$ is a certain graph parameter which can be as large as $\Omega(n)$ [Algorithmica '12]. The same relaxation can also be used to get a $k$-approximation algorithm, where $k$ is the number of different interfaces. This generalizes a similar result for the homogeneous cost case, where the cost of an interface is the same for all vertices.
Our main result is an $O(\log^2 n)$-approximation algorithm for \connectivity, which is the first non-trivial approximation for this problem. The algorithm is based on a similar LP relaxation with additional cut constraints to ensure connectivity. The rounding procedure resembles the one for \coverage but requires a more careful analysis to ensure that the connectivity constraints are satisfied.
\end{abstract}


\section{Introduction}
Designing Internet of Things (IoT) networks often demands establishing connection between large number of devices, which can be achieved by activating communication interfaces at each of them. However, the devices present in the network may vary according to the type of interfaces they have at their availability and the amount of resources a given interface consumes. For example some device may be able to utilize Bluetooth and Wi-Fi connections, while the others may provide Zigbee and Z-Wave interfaces. Since IoT devices are typically resource-constrained, it may be desirable to try to establish every relevant connection, while minimizing the amount of resources required to do so. 

\begin{figure}[t]
\centering
\input{figures/combined-illustration.tex}
\caption{From left to right: a common input graph, a covering assignment, and a connecting assignment for the same instance.}
\label{fig:example_problems}
\end{figure}
We model the above problem as a graph problem. Every device in the network is represented by a vertex with a set of labels encoding its available interfaces. Two devices can communicate if they share a common available interface and are connected by an edge, representing that they are within communication range. A connection is established between two adjacent devices whenever they share a common active interface. 
We want to find an assignment of active interfaces to vertices so that we ensure either the \coverage or \connectivity property. 

In the \coverage problem, we wish to establish every possible direct connection in the network which means that both endpoints of every edge must share a common activated interface (see the middle panel in \Cref{fig:example_problems} for an example) while in the \connectivity problem, we require all vertices to be (not necessarily directly) connected in the subgraph induced by active interfaces (see the right panel in \Cref{fig:example_problems} for an example). In both problems, we assume that activating an interface induces a cost that depends on the type of interface and on the vertex that activates it. The goal is to minimize the max cost, i.~e., the maximum cost of  interfaces activated at a given vertex across all vertices. This cost criterion is meant to spread the energy consumption somewhat evenly between devices, ensuring the sustainability of the whole network.

\subsection{Related Work}

In the following, let $k$ denote the number of different types of interfaces across the network, $n$ the number of vertices in the input graph, $m$ the number of edges, $\Delta$ the maximum degree, $c_{\max}$ the maximum cost and $c_{\min}$ the minimum cost of an interface. In the \emph{homogeneous} setting, the costs only depend on the interface while in the \emph{heterogeneous} case they depend on the interface and on the vertex. 

\subparagraph*{Coverage Problem.}
The algorithmic study of problems in multi-interface networks started with the problem of covering a multi-interface network, introduced in Caporuscio et al. \cite{EnergeticPerformanceOfServiceOrientedMultiRadioNetworks}. 
This problem has been mainly studied in the homogeneous setting, both for min-sum \cite{CostMinimizationInWirelessNetworks} and min-max \cite{MinimizeTheMaximumDutyInMultiInterfaceNetworks} objective functions. There has been recent work on the heterogeneous case for the problem of maximizing some profit function while keeping the energy cost within budget \cite{OnBudgetConstraineCoverageInMultiInterfaceNetworksBranchwidthAndTreewidthPerspectives}.
The min-sum variant is APX-hard for any constant $k \geq 3$ and unit costs, and it is hard to approximate within an $o\br{\ln k}$ factor even with unit costs, but it has a $\min \brc{k-1, \sqrt{n} \br{1+\ln n}}$-approximation \cite{CostMinimizationInWirelessNetworks}.
The min-max variant is \NPhard for any constant $\Delta \geq 5$, $k \geq 16$ and for unit costs, and it is not approximable within a $\eta \ln \Delta$ factor for some constant $\eta$ even on trees and with unit costs. On the positive side, using an approximation of \textsc{Set Cover}, it admits a $\br{\ln \Delta +1 +b \cdot \min \brc{c_{\max}, \ln \Delta +1}}$-approximation, where $b$ is a parameter that depends on the structural properties of the input graph and is bounded by many other graph structutal parameters such as treewidth, arboricity or maximum degree. It also has a $(1+\br{k-2} \frac{c_{\max}}{2 c_{\min}})$-approximation and a $(\frac{\Delta}{2} \frac{c_{\max}}{c_{\min}})$-approximation \cite{MinimizeTheMaximumDutyInMultiInterfaceNetworks}. 

\subparagraph*{Connectivity Problem.} The problem of connecting a multi-interface network was introduced by Kosowski et al. \cite{ExploitingMultiInterfaceNetworksConnectivityAndCheapestPaths}. It was first studied in the homogeneous setting, both for the min-sum \cite{ExploitingMultiInterfaceNetworksConnectivityAndCheapestPaths} and the min-max objective functions \cite{MinimizeTheMaximumDutyInMultiInterfaceNetworks}.
The min-sum variant where all costs are homogeneous is APX-hard for $k \geq 2$, $\Delta \geq 4$ and unit costs already, but has a $2$-approximation \cite{ExploitingMultiInterfaceNetworksConnectivityAndCheapestPaths}. This result was improved by Athanassopoulos et al. \cite{EnergyEfficientCommunicationInMultiInterfaceWirelessNetworks} to a randomized $(\frac{3}{2} + \epsilon)$-approximation for any constant $\epsilon$, using a reduction to a \textsc{Minimum Spanning Tree} problem on hypergraphs. Moreover, this approximation algorithm applies to the more general problem where for each interface, the set of edges that it can cover is specified in the input.

The min-max variant where all costs are homogeneous is \NPhard for any constant $\Delta \geq 3$, $k \geq 10$ and for unit costs. Similarly to the coverage, it is not approximable within a $\eta \ln n$ factor even for unit costs, but has simple $(1+(k-2) \frac{c_{\max}}{2 c_{\min}})$- and $\frac{\Delta}{2} \frac{c_{\max}}{c_{\min}}$-approximations \cite{MinimizeTheMaximumDutyInMultiInterfaceNetworks}.

Athanassopoulos et al. \cite{EnergyEfficientCommunicationInMultiInterfaceWirelessNetworks} expand the study of the min-sum variant to the case where only certain vertices, called \textit{terminals} need to be spanned and there may be several terminal groups. They provide a randomized $4$-approximation for this case. They also initiate the study of this problem in the heterogeneous setting. They show that for heterogeneous costs, the min-sum variant is hard to approximate within $o(\ln n)$ even when all vertices are terminals, but the min-sum group \textsc{connectivity} heterogeneous case admits a $O(\ln n)$-approximation.

\subparagraph*{Other Problems.} There has been some work on the study of the coverage problem from the perspective of parameterized complexity, possibly with an additional constraint on the maximum number of active interfaces in a vertex \cite{EnergyConsumptionBalancingInMultiInterfaceNetworks,MinMaxCoverageInMultiInterfaceNetworksPathwidth,OnComputationalComplexityOfCoveringMultiInterfaceNetworks}, yielding efficient algorithms in some graph classes. Some papers also study other problems in multi-interface networks, like cheapest path \cite{ExploitingMultiInterfaceNetworksConnectivityAndCheapestPaths}, matchings \cite{MaximumMatchingInMultiInterfaceNetworks}, or flow problems \cite{FlowProblemsInMultiInterfaceNetworks}.
Closely related combinatorial problems are \textsc{Set Cover}, \textsc{Minimum Spanning Tree} or \textsc{Steiner Tree Forest}, which are used for hardness results \cite{CostMinimizationInWirelessNetworks,MinimizeTheMaximumDutyInMultiInterfaceNetworks,OnComputationalComplexityOfCoveringMultiInterfaceNetworks}, and algorithms \cite{MinimizeTheMaximumDutyInMultiInterfaceNetworks,ExploitingMultiInterfaceNetworksConnectivityAndCheapestPaths,EnergyEfficientCommunicationInMultiInterfaceWirelessNetworks}.

\subsection{Our contribution and techniques}

We study the approximability of the \coverage and \connectivity problems minimizing the max cost under heterogeneous cost functions. We give a set-cover-style LP relaxation for the \coverage problem and show that treating the resulting fractional solution as a probability distribution and applying a randomized rounding procedure yields an $O\br{\log n}$-approximation with high probability. This is the first logarithmic approximation for this problem with heterogeneous costs, and it matches the previously known lower bound of $\Omega\br{\log n}$ for stars with unit costs \cite{MinimizeTheMaximumDutyInMultiInterfaceNetworks}. It also improves the previously best known $O\br{\ln \Delta + b \cdot \min\brc{\ln \Delta, c_{\max}}}$-approximation ratio \cite{MinimizeTheMaximumDutyInMultiInterfaceNetworks} (for some graph input parameter $b$), which was limited to the homogeneous case. 
A simpler deterministic rounding of the same LP gives a $k$-approximation, generalizing the $1+(k-2) \frac{c_{\max}}{2 c_{\min}}$-approximation for the min-max variant with homogeneous costs\cite{MinimizeTheMaximumDutyInMultiInterfaceNetworks}.

For the \connectivity problem, we extend the LP by adding flow-based cut constraints encoding global connectivity. Our rounding procedure is similar to the one for the \coverage problem; however, this time the process is iterative and the algorithm proceeds until all vertices are connected. Each iteration is randomized and we show that in expectation, in each such round the algorithm covers a significant fraction of edges of a graph that spans all vertices. This yields an $O\br{\log^2 n}$-approximation algorithm, which is the first non-trivial approximation algorithm for this problem.

\subsection{Paper Organization}

The rest of the paper is structured as follows:
In \Cref{sec:preliminaries}, we give notations, define the problems formally, and recall the elements of probability theory useful for the analysis of our algorithms. We also describe the preprocessing required for our randomized algorithms for the \coverage and \connectivity problems to work.
In \Cref{sec:coverage}, we present the $O\br{\log n}$-approximation and the $k$-approximation algorithms for the \coverage problem.
In \Cref{sec:connectivity}, we present the $O\br{\log^2 n}$-approximation algorithm for the \connectivity problem.
% We conclude with a short discussion on open questions and future work.
\section{Preliminaries} \label{sec:preliminaries}

Let $G=\br{V, E}$ be a connected simple graph. Let $\sbr{k}$ be the set of different types of interfaces across the network and let $\lambda\colon V \to 2^{\sbr{k}}$ be a function that assigns to each vertex $v\in V$ its set of available interfaces, such that for each edge $uv \in E$, $\lambda\br{u} \cap \lambda\br{v} \neq \emptyset$.
An assignment of active interfaces is a function $\lambda_A\colon V \to 2^{\sbr{k}}$ such that $\lambda_A\br{v} \subseteq \lambda\br{v}$ for every $v\in V$. We say that an edge $uv \in E$ is \emph{covered} if $\lambda_A\br{u} \cap \lambda_A\br{v} \neq \emptyset$. We denote by $\cov\br{\lambda_A} = \{uv \in E : \lambda_A\br{u} \cap \lambda_A\br{v} \neq \emptyset\}$ the set of edges covered by $\lambda_A$, and by $G_A$ the subgraph of $G$ with the vertex set $\bigcup_{e\in \cov\br{\lambda_A}}e$ and the edge set $\cov\br{\lambda_A}$.

In the \coverage problem, we wish to find an assignment of active interfaces $\lambda_A$ such that $G_A=G$. We call such an assignment a \emph{covering assignment}. In the \connectivity problem, we wish to find an assignment of active interfaces $\lambda_A$ such that, $V\br{G_A}=V$ and $G_A$ is connected. We call such an assignment a \emph{connecting assignment}, i.e., an assignment that induces a connected spanning subgraph of $G$.

We are also given a cost function $c\colon \sbr{k}\times V \to \mathbb{N}_{\geq 0}$, which denotes the cost of activating the interface $i\in\sbr{k}$ at the vertex $v\in V$. As a convention, we assume that $c(i,v)=\infty$ if $i\notin \lambda(v)$. For an assignment $\lambda_A$, define the \emph{cost} of $v$ as:
$$\COST\br{v, \lambda_A} = \sum_{i\in\lambda_A\br{v}}c\br{i, v}.$$
The \emph{max-cost} of an assignment $\lambda_A$ is defined as:
$$\COST \br{\lambda_A}=\max_{v\in V} \COST\br{v, \lambda_A}.$$
The goal of the min-max assignment problem is to find an assignment $\lambda_A$ that minimizes $\COST \br{\lambda_A}$ while ensuring either \coverage or \connectivity. 
The formal problem statements are as follows.

\pbdef{\coverage}{A connected simple graph $G=(V,E)$, available interfaces $\lambda\colon V \to 2^{\sbr{k}}$, a cost function $c\colon \sbr{k}\times V \to \mathbb{N}_{\geq 0}$.}{Find a covering assignment $\lambda_A$ minimizing $\COST\br{\lambda_A}$.}

\pbdef{\connectivity}{A connected simple graph $G=(V,E)$, available interfaces $\lambda\colon V \to 2^{\sbr{k}}$, a cost function $c\colon \sbr{k}\times V \to \mathbb{N}_{\geq 0}$.}{Find a connecting assignment $\lambda_A$ minimizing $\COST\br{\lambda_A}$.}

We denote by $\OPT$ the cost of an optimal solution to the problem at hand. We will refer to the min-max covering assignment problem as \coverage and the min-max connecting assignment problem as \connectivity. Note that there also exists a min-sum variant of these problems, where the goal is to minimize $\sum_{v\in V} \COST\br{v, \lambda_A}$ instead of $\COST\br{\lambda_A}$.
\subsection{Probabilistic tools}

Since our algorithms are randomized, we will need to use the following probability facts in the analysis. 
\begin{theorem}[Markov's Inequality]
Let $X$ be a non-negative random variable and $a > 0$. Then:
\[
    \Pr\sbr{X \geq a} \leq \frac{\mathbb{E}\sbr{X}}{a}.
\]
\end{theorem}

% \begin{theorem}[Wald's Identity \cite{OnCumulativeSumsOfRandomVariables}]
% Let $X_1, X_2, \dots$ be i.i.d.\ random variables with finite expectation $\mathbb{E}[X]$, and let $N$ be a non-negative integer-valued random variable independent of the sequence $(X_i)_{i\geq 1}$ with finite expectation $\mathbb{E}[N]$. Then:
% \[
%     \mathbb{E}\!\left[\sum_{i=1}^{N} X_i\right] = \mathbb{E}[N]\cdot\mathbb{E}[X].
% \]
% \end{theorem}

\begin{theorem}[(Weighted) Chernoff Bound]
    Let $X_1, X_2, \dots, X_n$ be independent random variables taking values in $\sbr{0,1}$, let $w_1, w_2, \dots, w_n$ be non-negative weights, such that for every $i\in \sbr{n}$ we have that $w_i\in \sbr{0,1}$. Let $X=\sum_{i=1}^{n} w_i\cdot X_i$ and let $\mu = \E{X}$. Then for every $\delta_1 > 0$ and $\delta_2 \in (0,1)$ we have that:
    \[
        \Pr\br{X \geq (1+\delta_1)\cdot \mu} \leq \br{\frac{e^{\delta_1}}{(1+\delta_1)^{1+\delta_1}}}^\mu
        \qquad
        \Pr\br{X \leq (1-\delta_2)\cdot \mu} \leq \br{\frac{e^{-\delta_2}}{(1-\delta_2)^{1-\delta_2}}}^\mu.
    \]
\end{theorem}

The proof of the above fact follows the proof of the standard Chernoff bound and we include it in \Cref{sec:appendix} for completeness.

\subsection{Preprocessing}
Before applying our algorithms (for both \coverage and \connectivity), we preprocess the instance in the following way: First, by a polynomial number of guesses and rescaling, we reduce the instances so that all of the costs are rational values in $\sbr{0,1}$ and we have $\OPT\geq 1/2$. Second, we partition vertices into two classes. A vertex $v$ is \emph{cheap} if $\sum_{i\in\lambda\br{v}} c\br{i,v}\leq 1$; for such vertices we may assume all available interfaces are activated. Vertices that are not cheap are \emph{expensive}; for them we may assume $\COST\br{v,\lambda_A}\geq 1$ in the analyzed solutions, while incurring only a factor $2$ loss in the approximation ratio. The full preprocessing procedure and its complete analysis are deferred to the Appendix \ref{sec:appendix-preprocessing}.
\section{Warmup: Approximation algorithms for the coverage problem} \label{sec:coverage}

In this section, we consider the \coverage problem with heterogeneous costs. According to the preprocessing, we assume that interface costs are in $\sbr{0,1}$ and $\OPT \geq 1/2$. 
The idea behind our ILP formulation and the algorithms is inspired by the rounding procedure known for \textsc{Set Cover} \cite{vazirani2001approximation} and related problems \cite{Angelidakis2018ShortestPQ}.


\subsection{ILP formulations of the \coverage problem}

We start with the ILP formulation of the \coverage problem. Let $x_{i, v}$ be a binary variable that is $1$ iff the interface $i\in\lambda\br{v}$ is assigned to vertex $v\in V$. If a given vertex $v$ is a cheap, then without large cost increase we can enforce $x_{i, v}=1$, for every $i\in\lambda\br{v}$. Otherwise, if a vertex $v$ is expensive, we can assume that in an optimal solution we have $\sum_{i\in \lambda\br{v}} x_{i, v}\cdot c\br{i, v} \geq 1$, which corresponds to the fact that $\COST\br{\lambda_A, v}\geq 1$. Thus, we can write the following ILP formulation of the \coverage problem:
\begin{align}
    \text{min} \quad & M \nonumber\\
    \text{s.t.} \quad & \sum_{i\in\lambda\br{v}}c\br{i, v}\cdot x_{i, v} \leq M \quad \forall v\in V \label{eqt:covering_cost}\\
    & \sum_{i\in\lambda\br{u}\cap\lambda\br{v}}\min\brc{x_{i, u}, x_{i, v}} \geq 1 \quad \forall uv\in E \label{eqt:coverage}\\
    &\begin{cases}
         x_{i, v} = 1 \quad \forall i\in\lambda\br{v} & \text{if } \sum_{i\in \lambda\br{v}} c\br{i, v} < 1\\
        \sum_{i\in \lambda\br{v}} c\br{i, v}\cdot x_{i, v} \geq 1 & \text{otherwise}
    \end{cases}\quad \forall v\in V \label{eqt:costs_coverage}\\
    & x_{i, v} \in \{0, 1\} \quad \forall v\in V, i\in\lambda\br{v} \label{eqt:integrality_coverage}
\end{align}
Constraints~\eqref{eqt:covering_cost} enforce that the cost of each vertex does not exceed the global cost $M$. Constraints~\eqref{eqt:coverage} ensure that for every edge $uv\in E\br{G}$, there exists a common activated interface. Constraints~\eqref{eqt:costs_coverage} encode the assumption that for every cheap vertex, all of its interfaces are activated, while for every expensive vertex, the cost of activated interfaces is at least $1$. Finally, constraints~\eqref{eqt:integrality_coverage} require the solution to be integral.
Note that the above formulation is not linear due to constraints \eqref{eqt:coverage}. However, by adding additional variables $z_{i, uv}$, we can linearize those by replacing them with the constraints $\sum_{i\in\lambda\br{u}\cap\lambda\br{v}}z_{i, uv} \geq 1$ together with $z_{i, uv} \leq x_{i, u}$ and $z_{i, uv} \leq x_{i, v}$ for every $i\in\lambda\br{u}\cap\lambda\br{v}$ and $uv\in E$. Thus, we can solve the LP relaxation of this formulation in polynomial time.

% This ILP formulation can easily be adapted to \textsc{min-sum} \coverage by removing \eqref{eqt:covering_cost} and replacing the objective function by $\sum_{v \in V} \sum_{i\in\lambda\br{v}} c\br{i, v} \cdot x_{i, v}$.


\subsection{A $k$-approximation algorithm for the Coverage problem}
In this section we show that a simple rounding of the LP solution can be used to obtain a $k$-approximation. Note that for this simple algorithm the preprocessing is not necessary, which makes it easily adaptable for the min-sum variant of the \coverage problem (assuming the ILP is appropriately changed).

\begin{algorithm}[H]
\caption{$k$-approximation algorithm} \label{alg:2}%A relaxation-based $k$-approximation algorithm.
Solve the LP relaxation and obtain a fractional solution $\{x_{i, v}\}_{i\in\lambda(v), v\in V}$.

Set $\lambda_A(v) \gets \{i \in \lambda(v) \colon x_{i, v} \geq \frac{1}{k}\}$, for every $v \in V$.

\Return{$\{\lambda_A(v)\}_{v\in V}$.}
\end{algorithm}

\begin{theorem}
\Cref{alg:2} is a $k$-approximation algorithm for the \coverage problem.
\end{theorem}

\begin{proof}
Consider an edge $uv\in E$. By averaging, there exists an interface $i\in\lambda\br{u}\cap\lambda\br{v}$ such that $\min\brc{x_{i, u}, x_{i, v}} \geq \frac{1}{k}$. Thus, $i\in\lambda_A\br{u}\cap\lambda_A\br{v}$ so $\lambda_A$ is a covering assignment. Moreover, since $x_{i, v} \geq \frac{1}{k}$ for every $i\in\lambda_A\br{v}$, we have that $c\br{\lambda_A\br{v}} = \sum_{i\in\lambda_A\br{v}}c\br{i, v} \leq k \cdot \sum_{i\in\lambda_A\br{v}}c\br{i, v} \cdot x_{i, v}$. Thus, we have $\COST\br{\lambda_A} \leq k \cdot \max_{v\in V}\sum_{i\in\lambda\br{v}} c\br{i, v} \cdot x_{i, v} \leq k \cdot \OPT$.
\end{proof}


\subsection{An $O\br{\log n}$-approximation algorithm for the \coverage problem}
In order to extract an integral solution with quality independent of $k$, we employ a randomized procedure. For each interface type $i\in\sbr{k}$, we draw a uniformly random threshold $t_i \in \sbr{0, 1}$ and activate the interface at a vertex if the corresponding variable scaled up by a factor of $O\br{\log n}$ exceeds $t_i$. The procedure is formalized by the following pseudocode.

\begin{algorithm}[H]
\caption{$O(\log n)$-approximation algorithm for the \coverage problem} \label{alg:1}%A relaxation-based $O(\log n)$-approximation algorithm.

Solve the LP relaxation and obtain a fractional solution $\{x_{i, v}\}_{i\in\lambda(v), v\in V}$.

Pick independent uniformly random thresholds $t_i \in \sbr{0, 1}$, for each $i \in \sbr{k}$.

$\lambda_A(v) \gets \{i \in \lambda(v) \colon 2\ln m\cdot x_{i, v} \geq t_i\}$, for every $v \in V$.

\Return{$\{\lambda_A(v)\}_{v\in V}$.}
\end{algorithm}

\begin{theorem}
\Cref{alg:1} is an $O(\log n)$-approximation algorithm for the \coverage problem which succeeds with high probability.
\end{theorem}
\begin{proof}
Firstly, we claim that the algorithm indeed returns a feasible solution:
    \begin{lemma}
With high probability, \Cref{alg:1} returns an assignment $\lambda_A$, such that for every edge $uv \in E$, $\lambda_A(u) \cap \lambda_A(v) \neq \emptyset$.
    \end{lemma}
\begin{proof}
    Consider an edge $uv\in E$. We have:
\begin{align*}
    \Pr\sbr{\lambda_A\br{u}\cap\lambda_A\br{v}=\emptyset} & = \prod_{i\in\lambda\br{u}\cap\lambda\br{v}}\Pr\sbr{t_i > 2\ln m \cdot \min\brc{x_{i, u}, x_{i, v}}}\\
    & = \prod_{i\in\lambda\br{u}\cap\lambda\br{v}}\br{1-2\ln m \cdot\min\brc{x_{i, u}, x_{i, v}}}\\&
    \leq \prod_{i\in\lambda\br{u}\cap\lambda\br{v}}\e{-2\ln m \cdot\min\brc{x_{i, u}, x_{i, v}}} 
    \\& = \e{-2\ln m \cdot\sum_{i\in\lambda\br{u}\cap\lambda\br{v}}\min\brc{x_{i, u}, x_{i, v}}}\\
    & \leq \e{-2\ln m} = \frac{1}{m^2}
\end{align*}

By applying the union bound over all $m$ edges we get that the algorithm returns a covering assignment with probability at least $1-\frac{m}{m^2} = 1-\frac{1}{m}$.
\end{proof}

Now, we also wish to show that the approximation factor achieved by our algorithm is upper bounded by $O\br{\log n}$. We have the following lemma:
\begin{lemma}
    \Cref{alg:1} returns a solution of cost at most $12\ln n \cdot \OPT$ with high probability.
\end{lemma}
\begin{proof}
Fix an expenssive vertex $v\in V$. Otherwise, if $v$ is cheap, then the cost of activating all interfaces at $v$ is at most $1\leq2\cdot\OPT$. We want to show that with high probability, the cost of $v$ is at most $12\ln n$ times the cost of $v$ in the fractional solution. 
By $X_{v, i}$ denote the random variable that is $1$ if $2\ln m\cdot x_{i, v} \geq t_i$ and $0$ otherwise. Then we can write $\COST\br{v, \lambda_A} = \sum_{i\in\lambda\br{v}}c\br{i, v} \cdot X_{v, i}$. We have: $\mathbb{E}\sbr{X_{v, i}} = \Pr\sbr{i \in \lambda_A\br{v}} = 2\ln m\cdot x_{i, v}$. Thus, $\mathbb{E}\sbr{\COST\br{v, \lambda_A}} = 2\ln m \cdot \sum_{i\in\lambda\br{v}}c\br{i, v} \cdot x_{i, v}$. Applying the Chernoff bound with $\delta = 2$ we get 
\begin{align*}
    \Pr\sbr{\COST\br{v, \lambda_A} \geq 3\cdot \mathbb{E}\sbr{\COST\br{v, \lambda_A}}} & 
    \leq \br{\frac{e^2}{27}}^{\mathbb{E}\sbr{\COST\br{v, \lambda_A}}}
    \leq \e{-\mathbb{E}\sbr{\COST\br{v, \lambda_A}}}\leq \frac{1}{m^2}
\end{align*}
where the last inequality is by using the fact that, $\sum_{i\in\lambda\br{v}}c\br{i, v} \cdot x_{i, v} \geq 1$. By applying the union bound over all $n$ vertices we get that with probability at least $1-\frac{n}{m^2}\geq 1-\frac{m+1}{m^2}$ for every expensive $v\in V$, $\COST\br{v, \lambda_A} \leq 6\ln m \cdot \sum_{i\in\lambda\br{v}}c\br{i, v} \cdot x_{i, v}$. Thus, with high probability we have $\COST\br{\lambda_A} \leq 6\ln m \cdot \max_{v\in V}\sum_{i\in\lambda\br{v}}c\br{i, v} \cdot x_{i, v} \leq 6\ln m \cdot \OPT\leq 12\ln n \cdot \OPT$. 
\end{proof}
Combining the above lemmas gives the desired result.
\end{proof}

\section{Approximation algorithm for the \connectivity problem} \label{sec:connectivity}

In this section we move towards the \connectivity problem with heterogeneous costs. We again assume that the instance is preprocessed, so we have that all costs are in $\sbr{0,1}$ and $\OPT \geq 1/2$.

\subsection{ILP formulation of the \connectivity problem}

We use a similar ILP to model the \connectivity problem, however we use additional flow variables to encode the connectivity requirement. The $x_{i, v}$ variables are defined as before. For each edge $e$, we add a variable $y_{e}$ which is $1$ iff the edge $e$ is covered by the solution. For any $S\subseteq V$, let $\delta(S)$ denote the set of edges with exactly one endpoint in $S$. Global connectivity is enforced via cut constraints: for every non-trivial cut $(S, V\setminus S)$, there needs to be at least one covered edge crossing the cut. Let $y\br{\delta\br{S}}=\sum_{e \in \delta(S)} y_e$. We can formulate the \connectivity problem as follows:
\begin{align}
    \text{min} \quad & M \nonumber\\
    \text{s.t.} \quad & \sum_{i\in\lambda\br{v}}c\br{i, v} \cdot x_{i, v} \leq M \quad \forall v\in V \label{eqt:4}\\
    & \sum_{i\in\lambda\br{u}\cap\lambda\br{v}}\min\brc{x_{i, u}, x_{i, v}} \geq y_{uv} \quad \forall uv\in E\\
    &\begin{cases}
            x_{i, v} = 1 \quad \forall i\in\lambda\br{v} & \text{if } \sum_{i\in \lambda\br{v}} c\br{i, v} < 1\\
            \sum_{i\in \lambda\br{v}} c\br{i, v}\cdot x_{i, v} \geq 1 & \text{otherwise}
    \end{cases} \quad \forall v\in V \label{eqt:5}\\
    & y\br{\delta\br{S}}\geq 1 \quad \forall S\subset V\colon S\neq\emptyset,\, S\neq V \label{eqt:6}\\
    & x_{i, v}  \in \{0, 1\} \quad \forall v\in V, i\in\lambda\br{v} \nonumber\\
    & y_{uv} \in \{0, 1\} \quad \forall uv\in E \label{eqt:integrality_connectivity}
\end{align}

Observe that the above ILP has an exponential number of constraints, but we can solve its LP relaxation in polynomial time using the ellipsoid method and a separation oracle for the cut constraints, which can be implemented using a max-flow algorithm.

\subsection{An $O(\log^2 n)$-approximation algorithm for the \connectivity problem}
Obtaining an integral solution for the \connectivity problem requires a more careful rounding procedure than for the \coverage problem. We wish to ensure that at least one connected spanning subgraph is covered by the assignment returned by the algorithm. To do so, we start by randomly sampling edges according to the $y$-variables of the fractional solution, in a way so that with high probability, the sampled subgraph $H$ is connected. To achieve this goal, we set the probability of choosing $e$ to be $p_e=\min\{1, 5\ln m \cdot y_e\}$. Then, the goal of the algorithm becomes covering $H$. To do so, we repeat the following process: we draw uniformly random thresholds $t_i$ for each type of interface $i\in \sbr{k}$, and activate $i$ at any vertex for which $5\ln m \cdot x_{i,v}\geq t_i$. Intuitively, each such iteration costs $O\br{\log n\cdot \OPT}$ in expectation and covers $\Omega\br{1}$ fraction of edges of $H$. After repeating this for a fixed number of $T = O\br{\log n}$ rounds, it can be shown that the algorithm covers all edges of $H$ almost surely. The formal description of the algorithm is given by the following pseudocode.

\begin{algorithm}[H]
\caption{$O\br{\log^2n}$-approximation algorithm for the \connectivity problem} \label{alg:3} %A relaxation-based $O(\log^2 n)$-approximation algorithm.

Solve the LP relaxation and obtain a fractional solution $\{x_{i, v}\}_{i\in\lambda(v), v\in V}, \{y_{uv}\}_{uv\in E}$.

For every edge $e\in E$, $p_e \gets \min\{1, 5\ln m \cdot y_e\}$. 

With probability $p_e$ set $Y_e \gets 1$, otherwise set $Y_e \gets 0$. 

Let $H$ be the subgraph of $G$ induced by the edges $e$ such that $Y_e = 1$.

$T \gets \frac{2\ln m}{1-1/e}$.

$\lambda_A\gets \{\emptyset\}_{v\in V}$.

\For{$j \in \sbr{T}$}{
Pick independent uniformly random thresholds $t_i \in (0, 1)$, for each $i \in \sbr{k}$.

$\lambda_A^j(v) \gets \{i \in \lambda(v) \colon 5\ln m \cdot x_{i, v} \geq t_i\}$, for every $v \in V$.

$\lambda_A \gets \lambda_A^j \cup \lambda_A$.
}
\Return{$\lambda_A$}
\end{algorithm}

\begin{theorem}
    \Cref{alg:3} is an $O\br{\log^2 n}$-approximation algorithm for the \connectivity problem which succeeds with high probability.
\end{theorem}
\begin{proof}
    In order to show that the algorithm returns a connecting assignment, we first show that with high probability, the subgraph $H$ is connected. It should be remarked that constructing $H$ explicitly is not necessary for the algorithm to work. Nonetheless, we do so for the sake of convenience of our analysis. We start with the following lemma:
    \begin{lemma}
        With high probability, $H$ is connected.
    \end{lemma}
    \begin{proof}
        Fix any $S\subset V$ with $S\neq\emptyset$ and $S\neq V$. By definition of the LP, there exists a natural number $\ell\in\mathbb{N}$, such that $\ell\leq y\br{\delta_H\br{S}}<\ell+1$. Let $\delta_H\br{S}=E\br{H}\cap \delta\br{S}$.
        We have:
        \begin{align*}
            \Pr\sbr{\spr{\delta_H\br{S}}=0} & = \prod_{e\in\delta\br{S}}\Pr\sbr{Y_e=0} = \prod_{e\in\delta\br{S}}(1-p_e) \\&\leq \prod_{e\in\delta\br{S}}\e{-p_e} \leq \e{-\sum_{e\in\delta(S)}\min\brc{1,\, 5\ln m\cdot y_e}}
            \\& \leq \e{-5\ell\ln m} = \frac{1}{m^{5\ell}}
        \end{align*}
        % where the last inequality uses the LP constraint $y\br{\delta\br{S}}\geq 1$ for every non-trivial cut, so $\sum_{e\in\delta(S)} y_e \geq \ell$.

        A cut in $G$ is called an $\alpha$-minimum cut if its value is less than $\alpha$ times the value of a minimum cut in $G$. 
        We now make use of the following Karger's lemma \cite{MinimumCutsInNearLinearTime, RandomSamplingInCutFlowAndNetworkDesignProblems}:
        \begin{lemma}[Karger's lemma]
            The number of $\alpha$-minimum cuts in a graph is at most $n^{2\alpha}$.
        \end{lemma}
        Observe that by construction of the LP, every non-trivial fractional cut has value at least $1$, so there are at most $n^{2\ell+2}$ such cuts that have fractional value less than $\ell+1$. Applying the union bound over all of them we obtain:
\begin{align*}
    \Pr\sbr{\exists\, S\subset V,\; S\neq\emptyset,\, S\neq V\colon \spr{\delta_H\br{S}}=0}
    &\leq \sum_{\ell=1}^{\infty}\sum_{\substack{S\subset V,\,S\neq\emptyset,\,S\neq V,\\ y\br{\delta_H\br{S}}<\ell+1}}\Pr\sbr{\spr{\delta_H\br{S}}=0} \\
    &\leq \sum_{\ell=1}^{\infty}n^{2\ell+2}\cdot \frac{1}{m^{5\ell}} \leq \sum_{\ell=1}^{\infty}\frac{1}{m^{\ell}} = \frac{1}{m-1}
\end{align*}
where the last inequality is due to the fact, that assuming that $G$ is not a tree, we get that $m\geq n$. Otherwise, \coverage and \connectivity are equivalent and one can use the $O\br{\log n}$-approximation algorithm instead.

Thus, with probability at least $1-\frac{1}{m-1}$, every non-trivial cut $S$ of $G$, has at least one edge of $H$ crossing it, which means that $H$ spans all vertices of $G$.
\end{proof}

Having shown that the subgraph $H$ is connected, we now wish to show that the assignment $\lambda_A$ returned by the algorithm covers all of the edges of $H$ almost surely. 
Let $j$ be an index of some iteration of the for loop. Let $H_j$ be the (possibly empty) subgraph of $H$ that is not covered by $\lambda_A$ at the beginning of the $j$-th iteration of the for loop. We wish to show, that in expectation, the solution $\lambda_A$ produced in the $j$-th iteration of the for loop, $\lambda_A^j$, covers a significant fraction of the edges in $H_j$. We have the following lemma:
\begin{lemma}
    We have
    $\mathbb{E}\sbr{\spr{\cov\br{\lambda_A^j, H_j}}} \geq \br{1-1/e} \cdot \mathbb{E}\sbr{\spr{E\br{H_j}}}$.
\end{lemma}
\begin{proof}
    Consider any edge $uv\in E\br{G}$. Observe that $\mathbb{E}\sbr{Y_{uv}}\leq 5\ln m \cdot y_{uv}$. By definition of the LP we get that:
    \begin{align*}
        \sum_{i\in\lambda\br{u}\cap\lambda\br{v}}\min\brc{x_{i, u}, x_{i, v}} & \geq y_{uv} \geq \frac{\mathbb{E}\sbr{Y_{uv}}}{5\ln m}\cdot
    \end{align*}
    Therefore, the probability that $uv$ is uncovered by $\lambda_A^j$ is at most:
    \begin{align*}
        \Pr\sbr{\lambda_A\br{u}^j\cap\lambda_A\br{v}^j=\emptyset} & = \prod_{i\in\lambda\br{u}\cap\lambda\br{v}}\Pr\sbr{t_i > 5\ln m\cdot\min\brc{x_{i, u}, x_{i, v}}} \\ 
        & = \prod_{i\in\lambda\br{u}\cap\lambda\br{v}}\br{1-5\ln m\cdot\min\brc{x_{i, u}, x_{i, v}}} \\& \leq \prod_{i\in\lambda\br{u}\cap\lambda\br{v}}\e{-5\ln m\cdot\min\brc{x_{i, u}, x_{i, v}}}\\ 
        & = \e{-5\ln m\cdot\sum_{i\in\lambda\br{u}\cap\lambda\br{v}}\min\brc{x_{i, u}, x_{i, v}}} \\& \leq \e{-5\ln m\cdot\frac{\mathbb{E}\sbr{Y_{uv}}}{5\ln m}}=\e{-\mathbb{E}\sbr{Y_{uv}}}
    \end{align*}
    
    Thus, we have $\Pr\sbr{\lambda_A\br{u}^j\cap\lambda_A\br{v}^j\neq\emptyset}\geq 1-\e{-\mathbb{E}\sbr{Y_{uv}} }$. By summing up over all edges in $H_j$ we lower bound the expected number of edges covered by $\lambda_A^j$ as follows:
    \begin{align*}
        \mathbb{E}\sbr{\spr{\cov\br{\lambda_A^j, H_j}}} & = \sum_{uv\in E\br{H_j}}\Pr\sbr{\lambda_A\br{u}^j\cap\lambda_A\br{v}^j\neq\emptyset}  
        \geq \sum_{uv\in E\br{H_j}}\br{1-\e{-\mathbb{E}\sbr{Y_{uv}} }}
        \\& \geq \br{1-1/e}\cdot \sum_{uv\in E\br{H_j}}{\mathbb{E}\sbr{Y_{uv}}} \geq \br{1-1/e} \cdot \mathbb{E}\sbr{\spr{E\br{H_j}}}
    \end{align*}
    where the second inequality is due to the fact that $1-e^{-x}\geq (1-1/e)\cdot x$, for $x\in\sbr{0,1}$ and the last inequality is obtained by linearity of expectation.
\end{proof}

Let $I$ be the random variable denoting the number of iterations needed to cover all edges of $H$, i.e., the smallest $j$ such that $\cov\br{\lambda_A, H} = E\br{H}$ after $j$ iterations (assume that the number of iterations might eventually get larger then $T$). We show that with high probability $I \leq T$, so the algorithm returns a feasible solution.
\begin{lemma}
    \label{lem:expected-iterations}
    $\Pr\sbr{I > T} = O\br{\frac{1}{m}}$.
\end{lemma}
\begin{proof}
    Let $\alpha = 1-1/e$. By definition, $\spr{E\br{H_{j+1}}}=\spr{E\br{H_{j}}}-\spr{\cov\br{\lambda_A^j, H_j}}$. By linearity of expectation and the previous lemma, $\mathbb{E}\sbr{\spr{E\br{H_{j+1}}}} \leq \br{1-\alpha}\cdot \mathbb{E}\sbr{\spr{E\br{H_{j}}}}$. Unfolding the recurrence gives $\mathbb{E}\sbr{\spr{E\br{H_{j}}}} \leq m e^{-\alpha(j-1)}$. By Markov's inequality, $\Pr\sbr{\spr{E\br{H_{j}}}\geq 1} \leq m e^{-\alpha(j-1)}$. Since $T = \frac{2\ln m}{1-1/e} = \frac{2\ln m}{\alpha}$, setting $j = T$ gives:
    $$\Pr\sbr{I > T} = \Pr\sbr{\spr{E\br{H_{T}}}\geq 1} \leq m e^{-\alpha(T-1)} \leq m e^{-2\ln m+\alpha} = O\br{\frac{1}{m}}.$$
    since $e^{\alpha}=O\br{1}$.
\end{proof}

It remains to show that the cost of the solution returned by the algorithm is at most $O\br{\log^2 n}\cdot \OPT$. We have with the following lemma:
\begin{lemma}
    With high probability, for every vertex $v\in V\br{G}$ and every iteration $j \in \sbr{T}$, $\COST\br{v, \lambda_A^j} \leq 15\ln m \cdot \OPT$.
\end{lemma}
\begin{proof}
Fix a vertex $v\in V$ and an iteration $j\in\sbr{T}$. If $v$ is cheap, then its overall cost is at most $1\leq 2\cdot\OPT$ regardless of activated interfaces. Thus, assume $v$ to be expensive, so $\sum_{i\in\lambda\br{v}}c\br{i,v} \cdot x_{i,v} \geq 1$.
By $X_{v,i}^j$ denote the indicator that is true iff $i\in \lambda_A^j$. Then, $\COST\br{v, \lambda_A^j} = \sum_{i\in\lambda\br{v}}c\br{i,v} \cdot X_{v,i}^j $. We have $\mathbb{E}\sbr{X_{v,i}^j} = 5\ln m\cdot x_{i,v}$, hence $\mathbb{E}\sbr{\COST\br{v,\lambda_A^j}} = 5\ln m\cdot \sum_{i\in\lambda\br{v}}c\br{i,v} \cdot x_{i,v}$. Applying the Chernoff bound with $\delta=2$ gives:
\begin{align*}
    \Pr\sbr{\COST\br{v, \lambda_A^j} \geq 3\cdot\mathbb{E}\sbr{\COST\br{v,\lambda_A^j}}}
    \leq \br{\frac{e^2}{27}}^{\mathbb{E}\sbr{\COST\br{v,\lambda_A^j}}}
    \leq e^{-\mathbb{E}\sbr{\COST\br{v,\lambda_A^j}}}
    \leq \frac{1}{m^5},
\end{align*}
where the last inequality uses $\mathbb{E}\sbr{\COST\br{v,\lambda_A^j}} = 5\ln m\cdot\sum_{i\in\lambda\br{v}}c\br{i,v} \cdot x_{i,v} \geq 5\ln m$. Hence $\COST\br{v,\lambda_A^j} \leq 15\ln m\cdot\sum_{i\in\lambda\br{v}} c\br{i,v} \cdot x_{i,v} \leq 15\ln m\cdot\OPT$ with probability at least $1-\frac{1}{m^5}$.

Applying the union bound over all $n \leq m + 1$ vertices and all $T \leq \frac{2\ln m}{1-1/e}$ iterations, the above holds for all $v\in V$ and $j\in\sbr{T}$ with probability at least $1 - \frac{T\cdot \br{m+1}}{m^5} = 1 - O\br{\frac{1}{m^3}}$.
\end{proof}
Consequently, summing up over all $T$ iterations, with high probability:
$$\COST\br{\lambda_A} \leq \sum_{j=1}^{T}\max_{v\in V\br{G}}\brc{\COST\br{v, \lambda_A^j}} \leq T\cdot 15\ln m\cdot\OPT = \frac{30\ln^2 m}{1-1/e}\cdot\OPT\leq \frac{120\ln^2 n}{1-1/e}\cdot\OPT$$

Combining the above lemmas gives the desired result.
\end{proof}


% \section{Combinatorial $O\br{\Delta}$-approximation algorithms}
\subsection{$\Delta$-approximation for the min-sum covering problem}
\begin{algorithm}[H]
\caption{A simple $\Delta$-approximation algorithm for the min-sum covering problem.}
Initialize $\lambda_A(v) = \emptyset$ for every $v \in V$.

For every edge $uv\in E$, let $i_{uv}=\argmin_{j\in\lambda\br{u}\cap\lambda\br{v}}\brc{c\br{j, u}+c\br{j, v}}$. 
Set $\lambda_A\br{u} = \lambda_A\br{u} \cup \{i_{uv}\}$ and $\lambda_A\br{v} = \lambda_A\br{v} \cup \{i_{uv}\}$.

\Return{$\{\lambda_A(v)\}_{v\in V}$.}
\end{algorithm}
\begin{theorem}
    The above algorithm is a $\Delta$-approximation algorithm for the min-sum covering problem.
\end{theorem}
\begin{proof}
    Consider any optimal solution $\lambda_A^*$. For any edge $uv\in E$, let $i_{uv}^* \in \lambda_A^*\br{u}\cap\lambda_A^*\br{v}$ be any interface that covers $uv$ in the optimal solution.
    For every vertex $v\in V$ and edge $uv\in E$ incident to $v$, we say that $v$ dominates $uv$, if $c\br{i_{uv}^*, v} \geq c\br{i_{uv}^*, u}$ (with ties broken arbitrarily). Let $D\br{v}$ be the set of edges dominated by $v$. We have:
    $$
    \OPT \geq \sum_{v\in V}\max_{uv\in D\br{v}}\brc{c\br{i_{uv}^*, u}+c\br{i_{uv}^*, v}}.
    $$
    The cost of the solution built by the procedure is upper bounded by:
    \begin{align*}
        \sum_{uv\in E}\min_{i\in \lambda\br{u}\cap\lambda\br{v}} \brc{c\br{i, u}+c\br{i, v}} &=\sum_{v\in V}\sum_{uv\in D\br{v}}\min_{i\in \lambda\br{u}\cap\lambda\br{v}} \brc{c\br{i, u}+c\br{i, v}}
         \\&\leq 
        \sum_{v\in V}\sum_{uv\in D\br{v}} \max_{uv\in D\br{v}}\brc{c\br{i_{uv}^*, u}+c\br{i_{uv}^*, v}} 
        \\&= 
        \sum_{v\in V}\Delta\cdot \max_{uv\in D\br{v}}\brc{c\br{i_{uv}^*, u}+c\br{i_{uv}^*, v}}
        \leq 
        \Delta \cdot \OPT
    \end{align*}
    Since by construction, the solution returned by the procedure is a covering assignment, the claim follows.
\end{proof}
\subsection{$1.39\Delta$-approximation for the min-sum \textsc{Connectivity} problem}
\begin{algorithm}[H]
\caption{A $1.39\Delta$-approximation algorithm for the min-sum \textsc{Connectivity} problem.}
Initialize $\lambda_A(v) = \emptyset$ for every $v \in V$.

For every edge $uv\in E$, let $i_{uv}=\argmin_{j\in\lambda\br{u}\cap\lambda\br{v}}\brc{c\br{j, u}+c\br{j, v}}$ and set $c_{uv} = c\br{i_{uv}, u}+c\br{i_{uv}, v}$. 

Let $T$ be a $1.39$-approximate Steiner tree of $G$ with terminal set $U$ and edge weights $c_{uv}$, computed using the algorithm of \cite{ByrkaGRS10}.

For every edge $uv\in T$, set $\lambda_A\br{u} = \lambda_A\br{u} \cup \{i_{uv}\}$ and $\lambda_A\br{v} = \lambda_A\br{v} \cup \{i_{uv}\}$.

\Return{$\{\lambda_A(v)\}_{v\in V}$.}
\end{algorithm}
\begin{theorem}
    The above algorithm is a $1.39\Delta$-approximation algorithm for the min-sum \textsc{Connectivity} problem.
\end{theorem}
\begin{proof}
    Consider any optimal solution $\lambda_A^*$ and let $G^*$ be the subgraph of $G$ induced by the edges covered in the optimal solution. Since $\lambda_A^*$ is a connecting assignment, $G^*$ connects all terminals in $U$. Let $T^*$ be a minimum Steiner tree of $G^*$ with terminal set $U$. For any edge $uv\in E\br{T^*}$, let $i_{uv}^* \in \lambda_A^*\br{u}\cap\lambda_A^*\br{v}$ be any interface that covers $uv$ in $T^*$.
    For every vertex $v\in V$ and edge $uv\in E\br{T^*}$ incident to $v$, we say that $v$ dominates $uv$, if $c\br{i_{uv}^*, v} \geq c\br{i_{uv}^*, u}$ (with ties broken arbitrarily). Let $D\br{v}$ be the set of edges dominated by $v$ in $T^*$. We have:
    $$
    \OPT\br{G} \geq \COST_A\br{T^*} \geq \sum_{v\in V}\max_{uv\in D\br{v}}\brc{c\br{i_{uv}^*, u}+c\br{i_{uv}^*, v}}.
    $$
    Note that for every $uv\in E\br{T^*}$ we have $c_{uv} \leq c\br{i_{uv}^*, u}+c\br{i_{uv}^*, v}$, so $\sum_{uv\in E\br{T^*}} c_{uv} \leq \COST_A\br{T^*} \leq \OPT\br{G}$. By the $1.39$-approximation guarantee of \cite{ByrkaGRS10}, the Steiner tree $T$ satisfies $\sum_{uv\in E\br{T}} c_{uv} \leq 1.39 \cdot \sum_{uv\in E\br{T^*}} c_{uv}$. The cost of the solution is therefore at most:
    \begin{align*}
        \sum_{uv\in E\br{T}}\min_{i\in \lambda\br{u}\cap\lambda\br{v}} \brc{c\br{i, u}+c\br{i, v}} &\leq 1.39\cdot\sum_{uv\in E\br{T^*}}\min_{i\in \lambda\br{u}\cap\lambda\br{v}} \brc{c\br{i, u}+c\br{i, v}}
        \\&=1.39\cdot\sum_{v\in V}\sum_{uv\in D\br{v}}\min_{i\in \lambda\br{u}\cap\lambda\br{v}} \brc{c\br{i, u}+c\br{i, v}}
         \\&\leq 
        1.39\cdot\sum_{v\in V}\sum_{uv\in D\br{v}} \max_{uv\in D\br{v}}\brc{c\br{i_{uv}^*, u}+c\br{i_{uv}^*, v}} 
        \\&= 
        1.39\cdot\sum_{v\in V}\Delta\cdot \max_{uv\in D\br{v}}\brc{c\br{i_{uv}^*, u}+c\br{i_{uv}^*, v}}
        \\&\leq 
        1.39\Delta \cdot \OPT
    \end{align*}
    Since $T$ is a Steiner tree spanning $U$ and every edge of $T$ is covered by the assignment, the solution returned by the procedure is a connecting assignment. The claim follows.
\end{proof}

\subsection{$2\Delta$-approximation for the min-max covering problem}

We start with the following observation. Let $M$ be any matching in $G$. Then, we have $\OPT\geq \max_{uv\in M}\brc{\min_{i\in \lambda\br{u}\cap\lambda\br{v}} \brc{\max\brc{c\br{i, u},c\br{i, v}}}}$. This is due to the fact, that for any edge $uv\in M$ there needs to exist a common interface $i\in \lambda_A\br{u}\cap\lambda_A\br{v}$ in the optimal solution. Since the above expression represents choosing the (localy) best interface for each edge in the matching, it is upper bounded by the cost of the optimal solution. We use the fact, that for any graph, there always exists a matching incident to every vertex of degree $\Delta$, except at most one. We also note that if $G'$ is a subgraph of $G$, then $\OPT\br{G'} \leq \OPT\br{G}$, since any covering assignment for $G$ is also a covering assignment for $G'$.
\begin{algorithm}[H]
\caption{$2\Delta$-approximation algorithm for the min-max covering problem}
If $n=1$ then return $\{\lambda\br{v}=\emptyset\}$, where $v$ is the only vertex in $G$.

Let $M$ be a matching in $G$, which is incident to the largest number of vertices of degree $\Delta$. Let $e=$ be an arbitrary edge incident to the remaining vertex of degree $\Delta$, if it exists. 

For every $H\in G-\br{M\cup\{e\}}$, call the procedure recursively and obtain a covering assignment $\lambda_A$ for all vertices of $H$.

For every edge $uv\in M\cup\{e\}$, let $i_{uv}=\argmin_{i\in\lambda\br{u}\cap\lambda\br{v}}\brc{\max\brc{c\br{i, u}, c\br{i, v}}}$. Set $\lambda_A\br{u} = \lambda_A\br{u} \cup \{i_{uv}\}$ and $\lambda_A\br{v} = \lambda_A\br{v} \cup \{i_{uv}\}$.

\Return{$\{\lambda_A(v)\}_{v\in V}$.}
\end{algorithm}
\begin{theorem}
    The above algorithm is a $2\Delta$-approximation algorithm for the min-max covering assignment problem.
\end{theorem}
\begin{proof}
    By construction, the solution returned by the procedure is a covering assignment. We show that its cost is at most $2\Delta\cdot \OPT$. We claim by induction that this is the case. For $n=1$, i. e., $\Delta=1$, the cost is $0$ and the claim holds. Assume, that the claim holds for some value of $\Delta$.
    By definition of $M$ we get that the cost of the solution is upper bounded by:
\begin{align*}
    &\max_{uv\in M\cup\brc{e}}\brc{\min_{i\in\lambda\br{u}\cap\lambda\br{v}}\brc{\max\brc{c\br{i, u}, c\br{i, v}}}} + \max_{H\in G-\br{M\cup\{e\}}}\COST_W\br{H} \\&\leq 
    2\cdot\OPT\br{G}+ \max_{H\in G-\br{M\cup\{e\}}} \brc{2\cdot\br{\Delta-1}\cdot \OPT\br{H}} \leq 2\Delta\cdot \OPT\br{G}
\end{align*}
where the first inequality is by the fact that $M$ and $\brc{e}$ are both matchings and the induction hypothesis.
\end{proof}


% \section{Approximation hardness of \textsc{Connectivity}}
In this section, we show that \textsc{Connectivity} cannot be approximated within a factor of $O\br{\log^{2-\epsilon} n}$ for any $\epsilon > 0$, unless $\text{NP} \subseteq \text{ZTIME}(n^{\text{polylog}(n)})$. To achieve this result we give a reduction from the \textsc{Group Steiner Tree} (\textsc{GST}) problem, which is known to be hard to approximate within a factor of $O\br{\log^{2-\epsilon} n}$~\cite{PolylogarithmicInapproximability}. In \textsc{GST}, we are given a graph $G = (V, E)$, a cost function $c: E \to \mathbb{R}_{\geq 0}$, and a collection of groups $g_1, \dots, g_k \subseteq V$. The goal is to find a subtree $T$ of $G$ that contains at least one vertex from each group minimizing $c(T)=\sum_{e\in E\br{T}}c\br{e}$. 

In order to prove this claim we need to use the specific properties of the construction used in order to show the hardness of approximation of \textsc{GST} given in~\cite{PolylogarithmicInapproximability}. 
Let $m$ be the size of some PCP system for $\text{NP}$. The instance of \textsc{GST} obtained in the reduction from this PCP system is a tree with the following properties:
\begin{enumerate}
    \item $k=m^{O\br{\log m}}$.
    \item The tree is a complete $d$-ary tree of height $H=O\br{\log^{\beta} m}=O\br{\log^{\frac{\beta}{\beta+2}} k}$ for some $\beta > 0$.
    \item The number of vertices of the tree is $O\br{d^H} = m^{O\br{\log^c m}}$ for some constant $c$.
    \item The cost of each edge is exactly $1/\tau$ times the cost of its parent edge for any desired constant $\tau > 1$.
\end{enumerate}
We call a tree of above type \emph{elegant}. We have:
\begin{theorem}[\cite{PolylogarithmicInapproximability}]
    \textsc{GST} cannot be approximated within a factor of $O\br{\log^{2-\epsilon} n}$ for any $\epsilon > 0$, unless $\text{NP} \subseteq \text{ZTIME}(n^{\text{polylog}(n)})$, even if the input graph is an elegant tree.
\end{theorem}
\begin{theorem}
    \textsc{Connectivity} cannot be approximated within a factor of $O\br{\log^{2-\epsilon} n}$ for any $\epsilon > 0$, unless $\text{NP} \subseteq \text{ZTIME}(n^{\text{polylog}(n)})$, even if the input graph is a DAG.
\end{theorem}
\begin{proof}
    The idea of our construction is as follows: The crucial difficuilty of the reduction is the fact that the cost criterion in \textsc{GST} is the sum of the costs of the edges, while in \textsc{Connectivity} it is the maximum cost of labels activated per vertex. In our construction, a desirable goal is for each edge outgoing of a vertex $v$ to represent a cost of all of the edges in the subtree rooted at the corresponding child of $v$. To achieve this, we will introduce a significant number of child representatives, each encoding a particular choice of cost of edges in the subtree rooted at $v$. In fact, each vertex will encode even richer information, i. e., the number of edges chosen at any level of the subtree rooted at $v$. Since the input tree is elegant, this information exatly determines the cost of the subtree rooted at $v$. Then, the next challenge is to force any reasonable algorithm to never exceed the cost entrusted to a given subtree by a too large factor. To achieve this, we will scale the costs in all subtrees using an appopriate (and usually differing) factor, so that each vertex in the tree has the same "allowed" cost. At last, we will also need to guess the value of the optimal solution to be able to scale the costs in the right way. Since, the number of possible values of the optimal solution is quasipolynomial this will not affect the running time.

    Assume, that we have guessed the value of the optimal solution $B$. Let $T$ be the input tree of \textsc{GST}. We construct a DAG, $D_B$ as follows: In the first stage of the construction  we assume $D_B$ is a rooted tree. Let $r$ be the root of $T$. Let $b_r=B$. We set $r'$ as the root of $D_B$ and we entrust it with a budget of $b_r$. The rest of the construction is recursive. The recurrence consists of a leaf $\ell$ of $D_B$ and a (some) root vertex $r$ of a subtree of $T$. For any child of $r$, call it $v$, we create $t_v=\spr{V\br{T_v}}^{O\br{H\br{T_v}}}$ vertex representatives $u_i$ of $v$, for $i\in [t_v]$ if the following conditions hols: For any such representative let $\brc{\mathcal{N}_{u_i}\br{j}}_{j=1}^{H\br{T_v}}$ be a sequence of non-negative integers such that $\mathcal{N}_{u_i}\br{j}\in \brc{0,\ldots, \spr{V\br{T_{v}}}}$ and there exists at least one subtree of $T_v$ with that number of vertices at each level. Note that this can be checked in quasipolynomial time, since the tree is elegant. Let $b_{u_i}=\sum_{j=1}^{H\br{T_v}}\mathcal{N}_{u_i}\br{j}\cdot \tau^j$ denote the budget entrusted to $u_i$. We demand that $b_{u_i}\leq b_{\ell}$. 
    
    Observe, that if an algorithm were to pick exactly the number of edges at each level of the subtree rooted at $v$ as encoded by $\mathcal{N}_{u_i}$, then the cost of the subtree rooted at $v$ would be exactly $b_{u_i}$. We connect the leaf $\ell$ to $u_i$ (with a directed edge) and create a label $u_i$ which we add to $\lambda\br{\ell}$ and $\lambda\br{u_i}$. We also set the cost of $u_i$ to be $b_{u_i}\cdot B/b_{\ell}$ for $\ell$ and $0$ for $u_i$. In this way we ensure that if the subtree $T_{\ell}$ is entrusted with a budget of $b_{\ell}$, then exceeding this budget $\alpha$ times in the original instance would correspong to paying $\alpha\cdot B$ in the new instance.
    
    In the second part of the construction, we need to encode the groups $g_1, \dots, g_k$. We assume that there is only one pair of source and sink subset in the \textsc{Connectivity} instance. The source conists solely of $r'$. For each group $g_i$ we create a vertex $s_i$ and we connect all the representatives of the vertices in $g_i$ to $s_i$ (directed). We also add terminal $g_i$ to both $\lambda\br{s_i}$ and all the representatives of the vertices in $g_i$. We set the cost of $g_i$ to be universally $0$. Finally, we set the sink subset to be $\brc{s_1, \dots, s_k}$.
    \begin{lemma}
        The size of $D_B$ is quasipolynomial in $m$.
    \end{lemma}
    \begin{proof}
        First, the number of sink nodes is $k = m^{O(\log m)}$. It remains to show that the number of representative vertices is quasipolynomial. Firstly, we show that each vertex has a quasipolynomial number of children. Let $r$ be such a vertex. Then, for each child vertex $v_i$, for $i\in[d]$, $r$ has $t_{v_i}=\spr{V\br{T_{v_i}}}^{O\br{H\br{T_{v_i}}}}$ children. Therefore the overall number of children is bounded by:
        \begin{align*}
            \sum_{i=1}^{d} t_{v_i} = \sum_{i=1}^{d} \spr{V\br{T_{v_i}}}^{O\br{H\br{T_{v_i}}}} \leq \sum_{i=1}^{d} \spr{V\br{T}}^{O\br{H\br{T}}} = d\cdot \spr{V\br{T}}^{O\br{H\br{T}}} = m^{O\br{\log^c m}}
        \end{align*}
        Since the depth of the DAG is at most the depth of the tree (plus $1$), we know that there exists an apropriate constant $c'$ such that the number of vertices in $D_B$ is at most \begin{align*}
            \sum_{i=0}^{H\br{T}+1} \br{\spr{V\br{T}}^{O\br{H\br{T}}}}^i = \sum_{i=0}^{H\br{T}+1} m^{O\br{\log^c m}\cdot i}  = \frac{m^{O\br{\log^c m}\cdot \br{H\br{T}+1}} - 1}{m^{O\br{\log^c m}} - 1} = m^{O\br{\log^{c'} m}}
        \end{align*}
        where the second inequality is by the fact that the sum is a geometric series and the last equality is by the fact that $H\br{T}=O\br{\log^{\beta}{m}}$.
    \end{proof}
    Now, we wish to show that upon guessing the value $B$, such that the optimal solution of the \textsc{GST} instance is $B$, any $\alpha$-approximation algorithm for \textsc{Connectivity} on $D_B$ would give an $\alpha$-approximation for the \textsc{GST} instance. 
        
    Firstly, without loss of generality we may assume that the algorithm picks at most one representative for each vertex of $T$. Indeed, if it picks more than one representative for a vertex $v$, then we can replace all the representatives of $v$ with a single representative $u$. Indeed, let $u_1,\dots, u_p$ be the representatives of $v$ picked by the algorithm. Let $\mathcal{N}_u\br{j}=\min\brc{d^j, \sum_{i=1}^{p} \mathcal{N}_{u_i}\br{j}}$ for each level $j$. It is easy to see that: Firstly, the cost of $u$ is at most the cost of all the representatives of $v$ picked by the algorithm. Secondly, the $\mathcal{N}_u\br{j}$ suffices to pick all edges in subtree rooted at $v$ that are picked by the representatives $u_1,\dots, u_p$. Therefore, we can replace all the representatives of $v$ with $u$ without increasing the cost of the solution and without losing any edge in the subtree rooted at $v$ that is picked by the algorithm.
        
    Now, we show how to recover the solution. We begin by removing all of the sink vertices from $D_B$ to obtain a tree. Initialize $\mathcal{T}$ to be an empty tree which will be the solution for \textsc{GST}. Then we traverse the tree in a top-down manner, starting from the root. Let $r$ be some vertex met along the way. Whenever an edge to a representative $u$ of a child of $r$, $v$ is picked, we append $uv$ to the solution of \textsc{GST} and we continue the traversal towards $u$. It is easy to see that the solution created in this way indeed contains at least one vertex from each group and thus is a feasible solution for \textsc{GST}. 
    \begin{lemma}
        If the optimal solution of the \textsc{GST} instance is $B$, then if there exists an $\alpha$-approximation algorithm for \textsc{Connectivity} on $D_B$, then the cost of $\mathcal{T}$ is at most $\alpha\cdot B$.
    \end{lemma}
    \begin{proof}
        We show this by a bottom-up induction over the tree $D_B$ (without the sink vertices). We claim, that 
        
        It is easy to see that the claim holds for the leaves of $D_B$. This is because each leaf $\ell$ of $D_B$ is a singular existing representative of some leaf of $T$ and thus, each edge of its parent 

    \end{proof}
\end{proof}


\bibliography{lipics-v2021-sample-article}

\appendix
\section{Weighted Chernoff Bound Derivation} \label{sec:appendix}


\begin{theorem}[(Weighted) Chernoff bound]
    Let $X_1, X_2, \dots, X_n$ be independent random variables taking values in $\sbr{0,1}$, let $w_1, w_2, \dots, w_n$ be non-negative weights, such that for every $i\in \sbr{n}$ we have that $w_i\in \sbr{0,1}$. Let $X=\sum_{i=1}^{n} w_i\cdot X_i$ and let $\mu = \E{X}$. Then for every $\delta_1 > 0$ and $\delta_2 \in (0,1)$ we have that:
    \[
        \Pr\sbr{X \geq (1+\delta_1)\cdot \mu} \leq \br{\frac{e^{\delta_1}}{(1+\delta_1)^{1+\delta_1}}}^\mu
        \qquad
        \Pr\sbr{X \leq (1-\delta_2)\cdot \mu} \leq \br{\frac{e^{-\delta_2}}{(1-\delta_2)^{1-\delta_2}}}^\mu.
    \]
\end{theorem}


\begin{proof}
We first prove the upper-tail bound for $\delta_1>0$. For any $\lambda>0$, Markov's inequality gives:
\[
\Pr\sbr{X \ge \br{1+\delta_1}\cdot\mu}
= \Pr\sbr{e^{\lambda X} \ge e^{\lambda\cdot\br{1+\delta_1}\cdot\mu}} \le \frac{\E{e^{\lambda X}}}{e^{\lambda\cdot\br{1+\delta_1}\cdot\mu}}.
\]
By independence, we have $
\E{e^{\lambda \cdot X}}
= \prod_{i=1}^n \E{e^{\lambda\cdot w_i\cdot X_i}}$.
Fix $i\in\sbr{n}$. Since $w_i\cdot X_i\in\sbr{0,1}$, by convexity of $y\mapsto e^{\lambda \cdot y}$ on $\sbr{0,1}$, for every $y\in\sbr{0,1}$,
$e^{\lambda \cdot y}
\le \br{1-y}\cdot e^0 + y\cdot e^{\lambda}
= 1 + \br{e^{\lambda}-1}\cdot y$.
Substituting $y=w_i\cdot X_i$, taking expectation, and using $1+t\le e^t$, we obtain:
\[
\E{e^{\lambda\cdot w_i\cdot X_i}}
\le 1 + \br{e^{\lambda}-1}\cdot w_i\cdot\E{X_i} \le \exp\br{\br{e^{\lambda}-1}\cdot w_i\cdot\E{X_i}}.
\]
Hence:
\[
\E{e^{\lambda X}}
\le \exp\br{\br{e^{\lambda}-1}\cdot\sum_{i=1}^n w_i\cdot\E{X_i}} = \exp\br{\br{e^{\lambda}-1}\cdot\mu}.
\]
Therefore,
$\Pr\sbr{X \ge \br{1+\delta_1}\cdot\mu}
\le \exp\br{\mu\cdot\br{e^{\lambda}-1-\lambda\cdot\br{1+\delta_1}}}$.
The function $f\br{\lambda}=e^{\lambda}-1-\lambda\cdot \br{1+\delta_1}$ is minimized at $\lambda=\ln\br{1+\delta_1}$, and so:
\[
\Pr\sbr{X \ge \br{1+\delta_1}\cdot\mu}
\le \exp\br{-\mu\cdot\br{\br{1+\delta_1}\cdot\ln\br{1+\delta_1}-\delta_1}} = \br{\frac{e^{\delta_1}}{\br{1+\delta_1}^{1+\delta_1}}}^{\mu}.
\]

Now let $\delta_2\in(0,1)$. For any $t>0$, again by Markov's inequality:
\[
\Pr\sbr{X \le \br{1-\delta_2}\cdot\mu}
= \Pr\sbr{e^{-tX} \ge e^{-t\cdot\br{1-\delta_2}\cdot\mu}} \le \frac{\E{e^{-t\cdot X}}}{e^{-t\cdot\br{1-\delta_2}\cdot\mu}}.
\]
Applying the same mgf estimate with $\lambda=-t<0$ yields $
\E{e^{-t\cdot X}}
\le \exp\br{\br{e^{-t}-1}\cdot\mu}$,
so:
\[
\Pr\sbr{X \le \br{1-\delta_2}\cdot\mu}
&\le \exp\br{\mu\cdot\br{e^{-t}-1+t\cdot\br{1-\delta_2}}}.
\]
The right-hand side is minimized at $t=-\ln\br{1-\delta_2}$, which gives:
\[
\Pr\sbr{X \le \br{1-\delta_2}\cdot\mu}
\le \exp\br{-\mu\cdot\br{\delta_2+\br{1-\delta_2}\cdot\ln\br{1-\delta_2}}} = \br{\frac{e^{-\delta_2}}{\br{1-\delta_2}^{1-\delta_2}}}^{\mu}.
\]
\end{proof}

\section{Preprocessing Details} \label{sec:appendix-preprocessing}

In this section, we provide the full preprocessing routine used for both the \coverage and \connectivity problems.

\subparagraph*{Scaling of the costs}
Firstly, we produce a polynomial number of instances, in such a way that for at least one of them we have a good lower bound on $\OPT$. For convenience, we allow costs to be rational numbers and scale them appropriately so that all of them lie in the interval $\sbr{0,1}$. To do so, by testing consecutive powers of $2$, we guess the smallest number $b$, such that the costliest activated interface among all vertices in the optimal solution has cost at most $2^b$. Let $C=\cl{\log\br{\max_{v\in V\br{G}, i\in \lambda\br{v}} \brc{c(i,v)}}}$. The number of the guesses is $C+1$, which is polynomial in the input size. If either some edge cannot be covered by interfaces of cost at most $2^b$, or there is no vertex $v$ and interface $i\in \lambda\br{v}$ such that $c\br{i,v}\in\sbr{1/2,1}$, then we discard such a guess. Otherwise, for each vertex $v\in V$, we temporarily discard from $\lambda\br{v}$ all interfaces with cost larger than $2^b$. Let $C_b=\max_{v\in V\br{G}, i\in \lambda\br{v}} \brc{c(i,v)}$ in this newly created instance. We normalize the costs of remaining interfaces by dividing them by $C_b$, so that the largest cost is exactly $1$. For the right guess of $b$, this gives $\OPT\geq C_b\geq 1/2$.

\subparagraph*{Cheap and expensive vertices}
Having established the scaling of costs, we now consider two types of vertices with respect to the new costs. For a vertex $v\in V$, if $\sum_{i\in \lambda\br{v}} c\br{i,v}\leq 1$, then we can safely assume that $v$ activates all interfaces in $\lambda\br{v}$, since this costs at most $1$ and is thus upper bounded by $2\cdot\OPT$. We call such vertices \emph{cheap}. The remaining vertices are \emph{expensive}. For every expensive vertex $v$, we can assume that for any output of the algorithm, $\COST\br{\lambda_A,v}\geq 1$, since enforcing this changes the cost by a factor of at most $2$ and does not affect asymptotic guarantees.

\subparagraph*{Finding the solution}
Since our algorithms are randomized, we also need to control the probability of failure. We ensure that, with high probability, for all values of $b$ we find a good solution. To do so, for each guess of $b$, we repeat the randomized procedure $K=\cl{\log_m\br{C}+1}$ times (which is polynomial). The algorithm is as follows.

\begin{algorithm}[H]
\caption{Preprocessing algorithm} \label{alg:preprocessing}
$\lambda_A(v) \gets \lambda(v)$ for every $v\in V$.

\For{$b\in \brc{0, \dots, C}$}
{
    For the remainder of this iteration, discard interfaces with cost larger than $2^b$ and divide all remaining costs by $\max_{v\in V\br{G}, i\in \lambda\br{v}} \brc{c(i,v)}$.

    \If{there is no solution for this new instance or there is no vertex $v$ and interface $i\in \lambda\br{v}$ such that $c\br{i, v}\in \sbr{1/2, 1}$}
    {
        \textbf{continue}
    }

    \For{$j\in \sbr{K}$}
    {
        $\lambda_A' \gets$ the output of the randomized procedure for the scaled-cost instance.

        \If{$\COST\br{\lambda_A'} < \COST\br{\lambda_A}$ and $\lambda_A'$ is a covering/connecting assignment}
        {
            $\lambda_A \gets \lambda_A'$.
        }
    }
}

\Return{$\{\lambda_A(v)\}_{v\in V}$.}
\end{algorithm}

Assuming that the failure probability of the randomized procedure is at most $O\br{1/m}$, for each fixed $b$ the probability of failure after $K$ repetitions is at most $O\br{1/m^K}=O\br{1/(mC)}$. Since there are $C$ guesses for $b$, a union bound implies that the probability that there exists some value of $b$ for which the algorithm fails is at most $O\br{1/m}$. In particular, this also holds for the correct guess of $b$.

\end{document}
